
Utilizando como base as métricas levantadas por \citeonline{Ferreira2020}, no
que tange ao consumo energético de nós sensores baseados em ESP32,
implementa-se a medição da corrente de pico durante o processamento e a
transmissão, da carga consumida por leitura e da autonomia energética do
dispositivo ao longo do tempo, métricas que pautam a validade da hipótese
\textit{H3}.

A corrente de pico ($I_{\text{pico}}$, em mA) é a maior corrente registrada
durante a execução da inferência embarcada ou do envio de pacotes via rádio
\gls{lora}. A autonomia energética estimada ($A$, em dias) de uma bateria de
capacidade nominal $C_{\text{bat}}$ (em mAh), submetida a uma corrente média
$I_{\text{médio}}$ (em mA) ao longo do tempo, é dada pela
Equação~\eqref{eq:autonomia_energetica}:

\begin{equation}
    \label{eq:autonomia_energetica}
    A = \frac{C_{\text{bat}}}{I_{\text{médio}} \times 24}
\end{equation}

A carga consumida em uma leitura é a integral da corrente $i(t)$ ao
longo de um ciclo completo, de um despertar ao seguinte, incluindo o repouso, calculada pela
regra do trapézio sobre as amostras e expressa em mAh:

\begin{equation}
    \label{eq:carga_leitura}
    Q_{\mathrm{leitura}} = \frac{1}{3600}\int_{t_{\mathrm{despertar},k}}^{t_{\mathrm{despertar},k+1}} i(t)\,dt
\end{equation}

\noindent com $i$ em mA e $t$ em segundos. Para estimar a autonomia em outros intervalos
entre leituras, o ciclo é modelado por dois estados: acordado, com corrente média
$I_{\mathrm{ac}}$ durante $t_{\mathrm{ac}}$, e em repouso, com corrente $I_{\mathrm{rep}}$
no restante do intervalo $T$. A autonomia $A$, em dias, para uma bateria de capacidade
$C_{\mathrm{bat}}$ em mAh, é então:

\begin{equation}
    \label{eq:autonomia_modelo}
    Q_{\mathrm{leitura}}(T) = \frac{I_{\mathrm{ac}}\,t_{\mathrm{ac}} + I_{\mathrm{rep}}\,(T - t_{\mathrm{ac}})}{3600}
    \qquad
    A = \frac{C_{\mathrm{bat}}}{Q_{\mathrm{leitura}}(T)} \cdot \frac{T}{86400}
\end{equation}

\noindent Essa expressão é a Equação~\eqref{eq:autonomia_energetica} aplicada à corrente
média do ciclo, $I_{\text{médio}} = 3600\,Q_{\mathrm{leitura}}(T)/T$. O modelo é ideal: não considera as perdas do regulador, a autodescarga nem a variação da
capacidade da bateria com a temperatura, de modo que a autonomia calculada é um limite
superior.